\section{Introduction}
\label{sec:intro}

LLM-based agents are increasingly deployed in enterprise workflows where they
read third-party content (emails, repository files, dispute evidence,
invoices) and call tools with real-world side effects \citep{yao2023react}.
Indirect prompt injection (IPI), ranked the top vulnerability for LLM
applications by OWASP \citep{owasp2025llmtop10}, occurs when adversarial content in
data the agent retrieves causes it to act contrary to the user's intent
\citep{greshake2023ipi}.

Existing IPI benchmarks have expanded from single-generation applications to
stateful and dynamic tool-using agents \citep{yi2024bipia, zhan2024injecagent, debenedetti2024agentdojo, li2026agentdyn}. Across these settings, however,
evaluated attacks remain predominantly attacker-authored instructions that
redirect the agent to an attacker-chosen task. Defenses
evaluated on these benchmarks therefore face primarily
instruction-bearing, redirecting attacks.

Instruction-like wording cannot itself define the security boundary. Useful
external content may contain instructions that an agent must follow,
especially when execution depends on information discovered at runtime
\citep{li2026agentdyn}. \citet{siu2026agentsecurity} frame instruction injection as
following a command from a source unauthorized to issue it. We argue that this
authority boundary should extend beyond commands: an
unauthorized source can steer a decision by asserting a false state, rule,
consequence, capability, or interpretation. External content constitutes IPI
when it attempts to influence a decision that its source is not
authorized to determine.
A false claim about a subject that the source is
authorized to determine is instead misinformation and ordinarily requires
independent evidence to detect. Neither linguistic form nor relevance to the
task establishes authority by itself. Because the boundary is set by the
decision, an attack also need not bring a task of its own: an environment
participant, such as a vendor or buyer, can tamper with the task in progress by
setting a wrong amount, recipient, or message in a call the task already
requires.

This distinction matters for defenses. Methods that mark external text or
train the model to separate instructions from data may still accept a
task-relevant claim as evidence \citep{hines2024spotlighting, wallace2024instruction, chen2024struq, chen2025secalign}, and injection
detectors may lose the instruction-like or off-task signals they rely on
\citep{li2025piguard, shi2025promptarmor}. Plan and policy defenses constrain calls
whose tools and arguments follow from the request, but a tampered call uses a
tool the task needs, and its consequential values are resolved at runtime from
retrieved content \citep{li2025drift, shi2025progent}, so it can pass
these checks. Evaluating only instruction-bearing, redirecting attacks can
therefore overstate robustness.

To measure this gap, we use environment-aware case construction. Each case
places attacker-controlled content that is relevant to the benign task in a
source the agent must process, identifies a decision that source is not
authorized to determine, and provides an agent-visible basis for resolving the
conflict. Holding the benign task, attacker role, targeted decision, and
harmful outcome fixed within each case, single-primitive attacks vary the
belief in the agent's decision process that the planted content directly
targets. This \emph{influence-primitive decomposition} provides a systematic
construction method for attacks beyond direct objective manipulation
(\S\ref{sec:framework}). We apply it to 90 validated cases across productivity,
DevOps, and marketplace environments, each with a compound attack and
single-primitive attacks for up to six primitives. Every case targets \emph{task
tampering}: the attack generally corrupts an action the task requires by
changing a structured field such as an amount or access level, or free-form
content such as an email body or resolution note.

\paragraph{Contributions.}

\begin{itemize}
\item
  \textbf{C1: An authority-based IPI boundary, operationalized.} We extend IPI's
  authority boundary from commands to any influence on a decision, and specify per
  case the attacker-controlled source, the decision outside its authority, and an agent-visible resolution basis.
\item
  \textbf{C2: A controlled task-tampering benchmark.} We co-design the benign task,
  attacker role, controlled source, and targeted decision of each case, and use
  influence primitives to vary how the planted content justifies the same
  harmful outcome.
\item
  \textbf{C3: A mechanism-level defense analysis.} We trace where an attack is
  stopped: before its content reaches the agent, by the agent itself, or at the
  tool call. This shows gaps in existing defenses at both ends: detectors miss
  non-instruction primitives more often, and action-level policies block few
  task-tampering calls.
\end{itemize}


